\section{{\fontsize{12pt}{14pt}\selectfont OBJETIVOS}}

\vspace{-0.3cm}
\subsection{{\fontsize{12pt}{14pt}\selectfont Objetivo geral}}
\novoBlocoTexto{
    Desenvolver e avaliar um protótipo web de um assistente educacional baseado em LLM e RAG, em que o RAG recupera trechos dos materiais didáticos de professores e a LLM, com base nesses trechos, classifica e organiza as dúvidas dos estudantes nos tópicos do plano de ensino, disponibilizando essas informações em um painel para acompanhamento dos conteúdos questionados pelos estudantes.
}
\subsection{{\fontsize{12pt}{14pt}\selectfont Objetivos específicos}}
\begin{itemize}
    \item Identificar, por meio da literatura, as dificuldades relacionadas à manifestação, recepção e organização das dúvidas discentes no contexto educacional;
    \item Realizar levantamento de requisitos funcionais e não funcionais do projeto;
    \item Modelar as interfaces de usuário;
    \item Desenvolver um protótipo web baseado em LLM e RAG que permita o registro das dúvidas discentes, a geração de respostas e a classificação dos questionamentos de acordo com os tópicos da disciplina;
    \item Implementar a geração de resumos individuais das dúvidas, e um painel docente para a visualização desses resumos organizados por tópicos com exibição da pergunta original dos estudantes e sem vínculo dos estudantes a estes dados;
    \item Avaliar, em ambiente controlado, o desempenho da classificação temática dasdúvidas realizada pela LLM, com apoio dos trechos recuperados pelo RAG, por meio da comparação com os tópicos de referência de um conjunto de dados público de perguntas de estudantes de Estrutura de Dados.
    
\end{itemize}
